\section*{Bab \ref{ch:g7:fractions} --- Pecahan: Membandingkan dan Menjumlahkan}

\begin{solution}{exo:g7:fractions:1}
$\dfrac{2}{5} = \dfrac{8}{20}$; \quad
$\dfrac{18}{24} = \dfrac{3}{4}$; \quad
$\dfrac{7}{3} = \dfrac{28}{12}$; \quad
$\dfrac{5}{9} = \dfrac{20}{36}$.
\end{solution}

\begin{solution}{exo:g7:fractions:2}
$\dfrac{12}{16} = \dfrac34$; \quad
$\dfrac{30}{45} = \dfrac23$; \quad
$\dfrac{27}{36} = \dfrac34$; \quad
$\dfrac{48}{12} = 4$.
\end{solution}

\begin{solution}{exo:g7:fractions:3}
$\dfrac{15}{35} = \dfrac37$ dan $\dfrac{21}{49} = \dfrac37$: senilai.

$\dfrac{16}{28} = \dfrac47$ dan $\dfrac{20}{36} = \dfrac59$: untuk
membandingkannya, $\frac47 = \frac{36}{63}$ dan
$\frac59 = \frac{35}{63}$ --- tidak senilai.
\end{solution}

\begin{solution}{exo:g7:fractions:4}
$\dfrac{7}{11} < \dfrac{9}{11}$ (\omterm{def:g4:fractions:def}{penyebutnya} sama).

$\dfrac56 = \dfrac{15}{18} > \dfrac{11}{18}$.

$\dfrac{13}{12} > 1$ sedangkan $\dfrac{19}{20} < 1$, jadi
$\dfrac{13}{12} > \dfrac{19}{20}$.
\end{solution}

\begin{solution}{exo:g7:fractions:5}
$\dfrac59 + \dfrac29 = \dfrac79$; \quad
$\dfrac{11}{7} - \dfrac47 = \dfrac77 = 1$; \quad
$\dfrac58 + \dfrac78 = \dfrac{12}{8} = \dfrac32$.
\end{solution}

\begin{solution}{exo:g7:fractions:6}
$\dfrac23 + \dfrac56 = \dfrac46 + \dfrac56 = \dfrac96 = \dfrac32$.

$\dfrac{7}{10} - \dfrac25 = \dfrac{7}{10} - \dfrac{4}{10} =
\dfrac{3}{10}$.

$\dfrac34 + \dfrac{5}{12} = \dfrac{9}{12} + \dfrac{5}{12} =
\dfrac{14}{12} = \dfrac76$.
\end{solution}

\begin{solution}{exo:g7:fractions:7}
$3 - \dfrac54 = \dfrac{12}{4} - \dfrac54 = \dfrac74$; \quad
$1 + \dfrac38 = \dfrac{11}{8}$; \quad
$2 - \dfrac76 = \dfrac{12}{6} - \dfrac76 = \dfrac56$.
\end{solution}

\begin{solution}{exo:g7:fractions:8}
$5 \times \dfrac{3}{10} = \dfrac{15}{10} = \dfrac32$; \quad
$\dfrac78 \times 4 = \dfrac{28}{8} = \dfrac72$; \quad
$\dfrac23$ dari $45$ adalah $\dfrac{2 \times 45}{3} = \dfrac{90}{3} = 30$.
\end{solution}

\begin{solution}{exo:g7:fractions:9}
Bersama-sama: $\dfrac14 + \dfrac38 = \dfrac28 + \dfrac38 = \dfrac58$
bagian pizzanya. Tersisa: $1 - \dfrac58 = \dfrac38$.
\end{solution}

\begin{solution}{exo:g7:fractions:10}
Yang dituangkan: $2 \times \dfrac16 = \dfrac26 = \dfrac13$ L. Tersisa:
\[
\frac34 - \frac13 = \frac{9}{12} - \frac{4}{12} = \frac{5}{12}
\text{ L}.
\]
\end{solution}

\begin{solution}{exo:g7:fractions:11}
Menjumlahkan besaran positif $\frac12$ pada $\frac35$ pasti memberi
hasil yang \emph{lebih besar} daripada $\frac35$. Tetapi
$\frac47 < \frac35$ (memang $\frac47 = \frac{20}{35}$ dan
$\frac35 = \frac{21}{35}$): \omterm{def:g1:addition:def}{jumlah} yang diakukan itu lebih kecil
daripada salah satu sukunya --- mustahil.
\end{solution}

\begin{solution}{exo:g7:fractions:12}
Selangkah demi selangkah:
\[
\frac12 + \frac14 = \frac34, \qquad
\frac34 + \frac18 = \frac68 + \frac18 = \frac78, \qquad
\frac78 + \frac1{16} = \frac{14}{16} + \frac1{16} = \frac{15}{16}.
\]
\omterm{def:g1:addition:def}{Jumlah} sebagiannya $\frac12, \frac34, \frac78, \frac{15}{16}, \dots$
setiap kali menjadi setengah kali lebih dekat ke $1$: pada setiap
langkah, bagian yang kurang terpotong menjadi setengahnya. \omterm{def:g1:addition:def}{Jumlahnya}
mendekati $1$ tanpa pernah mencapainya.
\end{solution}

\begin{solution}{pb:g7:fractions:1}
\textbf{1.} (a) $\frac12 + \frac14 + \frac14 = \frac24 + \frac14
+ \frac14 = \frac44 = 1$: sah. (b) $\frac14 + \frac18 +
\frac18 + \frac12 = \frac28 + \frac18 + \frac18 + \frac48 =
\frac88 = 1$: sah.

\textbf{2.} $\frac12 + \frac14 + \frac18 = \frac48 + \frac28 +
\frac18 = \frac78$: satu nada seperdelapan yang kurang.

\textbf{3.} Satu birama adalah $1 = \frac{16}{16}$: enam belas nada
seperenam belas. Seperdelapan adalah $\frac18 = \frac{2}{16}$: dua seperenam belas.

\textbf{4.} Setengah bertitik: $\frac12 + \frac14 = \frac34$.
\omterm{def:g3:division:half}{Seperempat} bertitik: $\frac14 + \frac18 = \frac28 + \frac18 = \frac38$.

\textbf{5.} Nada setengah bertitik bernilai $\frac34$ (pertanyaan 4):
tepat satu birama waltz. Isian yang lain, misalnya: \omterm{def:g3:division:half}{seperempat}
$+$ \omterm{def:g3:division:half}{seperempat} $+$ \omterm{def:g3:division:half}{seperempat} ($\frac14 + \frac14 + \frac14 =
\frac34$), atau \omterm{def:g3:division:half}{seperempat} $+$ seperdelapan $+$ seperdelapan $+$
\omterm{def:g3:division:half}{seperempat} ($\frac14 + \frac18 + \frac18 + \frac14 = \frac{2 + 1 + 1 +
2}{8} = \frac68 = \frac34$).

\textbf{6.} $\frac38 + \frac18 + \frac14 = \frac38 + \frac18 +
\frac28 = \frac68 = \frac34$: satu nada \omterm{def:g3:division:half}{seperempat} ($\frac14$) yang
kurang.

\textbf{7.} Seperdelapan bertitik: $\frac18 + \frac{1}{16} =
\frac{2}{16} + \frac{1}{16} = \frac{3}{16}$. \omterm{def:g3:division:half}{Seperempat}:
$\frac{4}{16}$. Yang \omterm{def:g3:division:half}{seperempat} lebih lama.

\textbf{8.} $\frac12 + \frac18 = \frac48 + \frac18 = \frac58$,
yaitu $\frac{10}{16}$: sepuluh seperenam belas.

\textbf{9.} $1 - \frac{11}{16} = \frac{16}{16} - \frac{11}{16}
= \frac{5}{16}$ bagian biramanya berupa kesenyapan.

\textbf{10.} Kedua contohnya berjumlah $\frac{2 + 1 + 1}{16} =
\frac{4}{16}$: sah. Semua irama yang mengisi empat seperenam belas
dengan potongan $2$ dan $1$:
\[
2{+}2, \quad 2{+}1{+}1, \quad 1{+}2{+}1, \quad 1{+}1{+}2, \quad
1{+}1{+}1{+}1 :
\]
ada lima irama.

\textbf{11.} Irama yang mengisi $n$ seperenam belas berakhir dengan
seperenam belas --- dan yang mendahuluinya mengisi $n - 1$ --- atau
dengan seperdelapan --- dan yang mendahuluinya mengisi $n - 2$.
Setiap irama terbilang tepat sekali dengan cara ini, jadi
$\text{cara}(n) = \text{cara}(n-1) + \text{cara}(n-2)$. Tabelnya:
\begin{center}
\renewcommand{\arraystretch}{1.2}
\begin{tabular}{c|cccccccc}
$n$ & $1$ & $2$ & $3$ & $4$ & $5$ & $6$ & $7$ & $8$ \\
\hline
cara & $1$ & $2$ & $3$ & $5$ & $8$ & $13$ & $21$ & $34$ \\
\end{tabular}
\end{center}
Satu nada setengah ($8$ seperenam belas) dapat ditabuh dengan $34$ cara.
($n = 4$ memperoleh kembali lima irama pada pertanyaan 10.)

\textbf{12.} Setiap bilangannya adalah \omterm{def:g1:addition:def}{jumlah} dua bilangan
sebelumnya: $3 + 5 = 8$, $5 + 8 = 13$, $8 + 13 = 21$,
$13 + 21 = 34$ --- aturan yang dibuktikan pertanyaan 11, berumur
seribu tahun dan kini dikenal sebagai aturan Fibonacci.

\textbf{13.} $\frac16 + \frac{1}{12} = \frac{2}{12} +
\frac{1}{12} = \frac{3}{12} = \frac14$: lama yang sama seperti
dua seperdelapan lurus ($\frac18 + \frac18 = \frac14$). Penukaran
swing itu sah.

\textbf{14.} Setiap nada triol berlangsung $\frac{1}{12}$: periksa,
$3 \times \frac{1}{12} = \frac{3}{12} = \frac14$
(\cref{prop:g7:fractions:mult}). Untuk nada setengah: masing-masing
berlangsung $\frac16$, karena $3 \times \frac16 = \frac36 = \frac12$.

\textbf{15.} \omterm{def:g3:division:half}{Seperempat} bertitik $+$ \omterm{def:g3:division:half}{seperempat} $+$ \omterm{def:g3:division:half}{seperempat}: dalam
seperdelapan, $3 + 2 + 2 = 7$ seperdelapan $= \frac78$: sah. Nada
setengah dan \omterm{def:g3:division:half}{seperempat} bernilai $4$ dan $2$ seperdelapan --- keduanya
bilangan \emph{\omterm{def:g3:numbers:evenodd}{genap}}. \omterm{def:g1:addition:def}{Jumlah} \omterm{def:g3:numbers:evenodd}{bilangan genap} mana pun tetap \omterm{def:g3:numbers:evenodd}{genap},
tetapi birama tujuh-delapan memerlukan $7$ seperdelapan, yaitu
\omterm{def:g3:numbers:evenodd}{bilangan ganjil}: mustahil. Untuk mengisi birama yang \omterm{def:g3:numbers:evenodd}{ganjil}, sesuatu
yang \omterm{def:g3:numbers:evenodd}{ganjil} --- seperdelapan atau nada bertitik --- harus muncul.

\end{solution}
